\section{Continuing an individual Gamma marking}\label{sec:markedpaths}

Let $\mathcal C_0$ denote the mirror category of Theorem~\ref{thm:diskgeneration}, with perfect envelopes understood, and write $D_s$ for its geometric disks. Theorem~\ref{thm:coherentcollection} gives an equivalence $U:\mathcal C_0\to\mathrm{Perf}(X)$ which describes their symmetries. We seek an ordinary equivalence on $\mathcal C_0^{\mathrm{op}}$ computing their individual classes. Taking the opposite reverses the order of the collection. As in Proposition~\ref{prop:ordinarymarking}, we will compare it with $D_X U$ on this same source.

\subsection{A regular ray potential}
For independent ray coefficients $c=(c_i)\in(\C^*)^m$ write $W_c=\sum_i c_i x^{v_i}$. There are two coefficient tuples of interest. The tuple $*=(t,\ldots,t)$ preserves the fan group and supplies the end of the compact perturbation. A second positive tuple $a$ lies in the actual small chart of \cite[\S3.1]{FZ} and has distinct Morse values. We explain both its existence and the connection between the charts.

The incidence $\sum c_i x^{v_i} v_i=0$ is smooth of dimension $m$ since the derivatives in the coefficient directions span $N_\C$. Its projection to coefficient space is finite. Indeed in a valuative boundary specialization with all $c_i$ units, a nonzero valuation vector for $x$ would give, at the terms of minimal valuation, a linear relation with nonzero coefficients among the vertices of a proper face. These vertices are independent. Applying this argument in a product of projective lines compactifying the torus proves properness. Each fibre is proper and affine, hence finite; it has the Newton length of Proposition~\ref{prop:cluster}. Smoothness of the incidence and regularity of the base give flatness by miracle flatness, and in characteristic zero there is a dense étale locus. On its off-diagonal space of pairs of sheets, equality of critical values cannot be identical on any component. Differentiating the equality would give $x^{v_i}=x'^{v_i}$ for every $i$ and thus $x=x'$. Removing the finite image of the coincidence locus gives a dense Zariski open of coefficients with distinct Morse values. It meets every positive real open box by the identity principle.

On the small slice the divisor coordinate is $\tau(c)=\sum_i D_i\log c_i$ at $Q=1$, using the real logarithms at positive coefficients. To see this in our chart put $y_i=c_i/t$. In the small toric $I$-function after removing the classical divisor exponential, the highest possible power of $z$ in curve degree $d\ne0$ is $-c_1\cdot d-\#\{i:D_i\cdot d<0\}$. A negative intersection contributes its factor at zero. If $c_1\cdot d=1$, all intersections cannot be nonnegative, since the sole nonzero intersection would be one and $\sum_i(D_i\cdot d)v_i=0$. Thus the series after removal has constant term one and no term at $z^{-1}$. In the normalized mirror factorization the volume column corresponds to the unit and the divisor shift is $\sum D_i\log y_i$, by \cite[Lemma~4.32 and (4.14)--(4.15)]{CCIT}. Restoring $c_1\log t$ gives $\tau(c)$. The fan sequence gives $q(c)^d=\prod_i c_i^{D_i\cdot d}$, and any two coefficient lifts of $q$ differ by a torus translation preserving the logarithmic volume. Use the global small mirror identification along the divisor torus \cite[\S7.3]{GHIKLS}. Choose the anchor inside both the positive small chart of \cite{FZ} and the analytic mirror neighbourhood of \cite[Proposition~6.15, 6.16(2), and (6.11)]{IriGlobal}. To see that these meet, in the nef coordinates of \cite[\S3.1]{FZ} take all $q_j=s>0$. The sum of the nef basis divisors is ample, so every nonzero effective monomial tends to zero as $s\downarrow0$. This approaches the large-radius point of both charts. The regular distinct-value locus descends to the coefficient-torus quotient by translation, and therefore meets a small positive box in this intersection. Choose a positive coefficient lift $a$ there. The finite Brieskorn lattice is locally free in the neighbourhood of the large-radius point just specified; smoothness supplies the lattice-basis condition since every maximal cone is unimodular. The angle of our initial ray can be chosen arbitrarily close to zero and admissible at both endpoints.

Write $h_\varphi(u)=\operatorname{Im}(e^{-i\varphi}u)$. In the part of the value disk containing the critical values, a preferred path from $u$ follows $u+\mathbb R_{\ge0} e^{i\varphi}$. Beyond these values bend the paths with the same lift of the direction to radial ends in a small positive angular interval adjoining the stopped page. At a tie use one common path and distinguish its disks by the critical points. The geometric exceptional order increases with $h_\varphi$, the order on the opposite category decreases.

\begin{lemma}\label{lem:regularpath}
There is a path from the chosen compact perturbation to the uncut $W_a$ with all its critical points Morse, all its critical values distinct after the initial time, and freely prescribed relative height order in the initial tied blocks. A distinguished system of initial common paths, split in that order and then carried by isotopy, gives an ordinary equivalence of the endpoint stopped categories carrying each initial disk to its dragged disk.
\end{lemma}

\begin{proof}
Let $F(y)$ be the finite Laurent deformation of Proposition~\ref{prop:cluster}, $\rho$ the compact cutoff of Section~\ref{sec:compactmirror}, and use
\[
                         W_c+\rho(F(y)-W_c).
\]
The two variables play different roles: the critical cluster is determined by $y$, the far fibres by $c$. In a small ball about the central $y^*$ the discriminant of the characteristic polynomial of multiplication by $F(y)$ on the retained finite flat algebra is holomorphic. It is nonzero at some positive ray tuple $b$ arbitrarily close to $*$ by the preceding calculation. The independent basis derivatives in the proof of Lemma~\ref{lem:quot} make the mirror map a submersion at the symmetric Morse point. The labelled canonical coordinates can therefore move with prescribed small transverse velocities, including at coincident values. Lift this splitting by a local section. Join its regular end to the point with $F(y)=W_b$ in the ball with the discriminant removed: on the complex line between the endpoints one can detour around the isolated zeros. Move $c$ from $*$ to $b$ in their small common neighbourhood. On the compact cutoff annulus the real and imaginary differentials retain positive Poisson bracket by $C^1$-closeness to the central map. Outside the cutoff the estimates of Lemma~\ref{lem:rayest}, uniform on a compact coefficient neighbourhood, exclude other critical points at large radius. On the intervening compact set outside the cluster the differentials remain nonzero by uniform closeness. At the end the function is precisely $W_b$. From $b$ to $a$ take an uncut coefficient path in a narrow tube about the straight segment in their real logarithms. Isolated zeros of the ray discriminant on the complexified line can again be bypassed. This choice fixes the lift of both real logarithms.

For the geometric transport, take $R$ larger than all the critical-value moduli and the images of the perturbation supports. Over the initial splitting translate the common arcs in parallel strips near their critical points, bending to radial arcs at the boundary as in Lemma~\ref{lem:samedisks}. Thereafter drag by a base isotopy translating small balls of the moving values and fixed near $|w|=R$. It is enough to consider finitely many smooth subintervals. In a Morse chart the unstable manifold for the gradient of $\operatorname{Re}(e^{-i\varphi}\pi)$ gives a smooth family of collapsing balls: its Hessian has $n$ positive and $n$ negative eigenvalues, and its imaginary part is constant on gradient lines. Compact local caps bounded at a small positive projected distance thus vary smoothly by the parameter-dependent unstable manifold theorem \cite[Lemmas~1.13--1.14]{Sei}. Take separate charts at the points of a tied fibre. The remaining compact arcs lift by parallel transport. On their regular compact subarcs the lifts are complete by the exponential gradient lower bound of Lemma~\ref{lem:rayest} outside a compact set, and vary smoothly with initial data. At a tie their disjoint spheres follow the same parallel transport. The resulting disks and their end collars have compact traces over each closed parameter interval.

Let the function at time $v$ be $\pi_v$. Truncate it by $|\pi_v|\le R$, $\Phi\le L$ and mark the page $\pi_v=R$, $\Phi\le L'<L$. Take $L'$ above the uniform regular-page threshold, then $L$ beyond all disk traces and end collars. For the primitive $\lambda_\Phi+\kappa\pi_v^*\alpha_0$ the fibre connection is unchanged. At large logarithmic radius $r$ in the bounded base, let $q$ be the squared conformal scale of $d\pi_v$ on its complex horizontal line. The horizontal projection of the Liouville field to the base is
\[
                  \frac{d\pi_v(Z_{\lambda_\Phi})+\kappa q\,w/2}
                       {1+\kappa q}.
\]
Here $q$ grows at least exponentially with $r$. The vertical derivative of $\Phi$ along its fibre gradient is bounded below by $c r^2$ on the large horizontal face by Lemma~\ref{lem:rayest}. The horizontal lift of a bounded base vector has size $O(q^{-1/2})$. Thus the horizontal face is outward uniformly for $\kappa\ge1$. On the vertical face the projection tends to $w/2$ at infinity. On its remaining compact part the positive base area $\kappa$ dominates the bounded native horizontal area and covector. Fixing $\kappa$ large gives the same strict outward inequality there. In dimension one the large horizontal face is empty. Round the seam away from the ribbon and terminal spheres. These estimates also apply on a small annulus about $|w|=R$.

On the collar of each disk subtract the differential of its intrinsic primitive, extended to a tubular neighbourhood with a cutoff. The base primitive vanishes on the radial end of the arc, so the required local correction is independent of $\kappa$ there. Increasing $\kappa$ preserves strict outwardness under this bounded correction. The primitive now annihilates the disk tangents on that collar; since the disk is Lagrangian its Liouville field is tangent there. This makes a smooth family of cylindrical decorated disks. Their links avoid the reference ribbon. The flattening in Proposition~\ref{prop:compactpair} may begin beyond all initial traces. Its local Kähler adjustments at the centres can be undone with the same cap argument. The old and new unflattened carriers have the same vertical germs on both the disks and the page, and complete to the same torus by the exhaustion argument of that proposition. This also holds after the common bounded end correction: on its compact annular support the radial speed is strictly positive throughout, so a backward trajectory spends uniformly bounded time in that annulus before the unchanged exhaustion argument applies. Enlarging a truncated page through regular collar levels is a deformation of embedded ribbons. These observations identify the incoming objects with those of Theorem~\ref{thm:diskgeneration}.

Under Moser identification the disks give an exact cylindrical isotopy, and the moving ribbon and their links remain disjoint. Choose the thickenings in Lemma~\ref{lem:geometrictransport} small enough to avoid all these links, also at nearby times by compactness. Over each small interval the isotopy identifies all corresponding branes in the common smaller-stop category \cite[\S2.4]{GPS}. Its forgetful arrows act by the same branes. Composition gives one equivalence with all the asserted images.
\end{proof}

\subsection{Which disk changes at a wall}
The dragged system at $a$ may differ from the preferred rays there. The following calculation removes any need to identify the two systems in advance. Write $\mathcal C_a$ for the category of the terminal compact pair of Lemma~\ref{lem:regularpath}.

\begin{proposition}\label{prop:markreduction}
Suppose a complex-linear equivalence $\mathcal C_a^{\mathrm{op}}\to\mathrm{Perf}(X)$ gives each preferred disk the unsigned Gamma class of its labelled critical point at $a$. Then there is an equivalence on $\mathcal C_0^{\mathrm{op}}$ giving each $D_s$ the unsigned Gamma class of $s$.
\end{proposition}

\begin{proof}
The local critical cluster has the canonical coordinates of the analytic big quantum product. Along the uncut divisor path there are ambient analytic big germs with the same property. For completeness fix nef divisor coordinates and a strictly positive $q_0$ in the initial domain of absolute convergence with positive radii for the remaining variables \cite{IriConv}. On our compact divisor path $|q^d|/q_0^d\le e^{A c_1\cdot d}$ for effective $d$, by ampleness and compactness of the unit $c_1$-slice in the effective cone. In a product coefficient with three fixed insertions of complex degrees $p_1,p_2,p_3$ and further insertions of degrees $p_j'$ and multiplicities $m_j$, the dimension condition reads
\[
         c_1\cdot d=p_1+p_2+p_3-n+\sum_j(p'_j-1)m_j .
\]
The divisor equation and this bound permit absolute majorization by the original series after shrinking the transverse radii by fixed factors. The germs thus agree on overlaps and give semisimple coordinate neighbourhoods along the regular path. Their calibrated small slice uses precisely the divisor logarithms above.

Keep the initial splitting and endpoint collars. Perturb the remaining labelled value path slightly so it crosses equal-height walls transversely at finitely many simple crossings. One may use a fine polygonal approximation avoiding codimension-two intersections, rounding away from the crossings. The local inverses of the canonical coordinates lift the perturbation to an endpoint-fixed homotopic path of analytic germs. Geometrically add the corresponding small constants to $\pi_v$ in disjoint moving Morse charts, cutting them off on compact regular annuli. The uniform positive bracket there persists. This achieves precisely the perturbed values, and Lemma~\ref{lem:regularpath} still transports the geometric disks. We use the analytic germs for continuation of periods throughout this argument.

Rotate the picture so the rays point right. At a crossing call the point on the left $r$ and the one on the right $f$. The ray of $f$ survives, as do all spectator rays. Only the new ray from $r$ bypasses $f$ on the other side. In the common comparison of Lemma~\ref{lem:regularpath} all unaffected disks are identified. On the post-crossing fibre one may end the dragged and newly preferred rear arcs on opposite sides of the common front end in the stop-free interval. The end gauges then have disjoint terminal collars. Flattening beyond the compact disk traces and undoing the local cap adjustment as in Lemma~\ref{lem:regularpath} identifies both systems in the same category of the compact fibration. Each is exceptional by the regular case of Proposition~\ref{prop:diskhoms}. Use the Euler form $B$ on the opposite category, with its decreasing-height order. Expand the new rear class $v$ in the incoming integral basis. If the adjacent slots are $i,i+1$, the vanishings $B(v,e_j)=0$ in the earlier slots eliminate the coefficients there successively from the left. The equations $B(e_j,v)=0$ in the later slots eliminate them from the right. Thus the only possibilities are
\[
\begin{array}{ll}
 (r,f)\ \longrightarrow\ (f,v),&v=\pm(r-B(r,f)f),\\
 (f,r)\ \longrightarrow\ (v,f),&v=\pm(r-B(f,r)f).
\end{array}
\]
Indeed orthogonality determines the line and both $v$ and the indicated vector have self-Euler pairing one. Each move preserves integrality and the basis property. The same computation can start at the terminal endpoint, whose preferred classes form a basis by hypothesis.

On the analytic side the front Laplace ray survives the same crossing. The clockwise-first pairing of Gamma sections is the coherent Euler form, and the two changes of the rear column are exactly the above elementary mutations \cite[\S\S2.5--2.6 and 4.3]{GGI}, \cite[Remark~4.13]{GI}. On a wall-free interval the columns are flat under the chosen continuation and the rays follow isotopy. Starting at the calibrated terminal end and traversing the walls in reverse therefore computes every labelled class up to sign. At the final coalescence the disjoint Morse caps approach the specified individual $D_s$. The asymptotic fundamental solution extends with its prescribed leading normalized idempotent even when exponents coincide \cite[Proposition~2.5.1]{GGI}. This identifies each label through the tie and proves the assertion.
\end{proof}
